% Options for packages loaded elsewhere
\PassOptionsToPackage{unicode}{hyperref}
\PassOptionsToPackage{hyphens}{url}
\documentclass[
  ignorenonframetext,
]{beamer}
\newif\ifbibliography
\usepackage{pgfpages}
\setbeamertemplate{caption}[numbered]
\setbeamertemplate{caption label separator}{: }
\setbeamercolor{caption name}{fg=normal text.fg}
\beamertemplatenavigationsymbolsempty
% remove section numbering
\setbeamertemplate{part page}{
  \centering
  \begin{beamercolorbox}[sep=16pt,center]{part title}
    \usebeamerfont{part title}\insertpart\par
  \end{beamercolorbox}
}
\setbeamertemplate{section page}{
  \centering
  \begin{beamercolorbox}[sep=12pt,center]{section title}
    \usebeamerfont{section title}\insertsection\par
  \end{beamercolorbox}
}
\setbeamertemplate{subsection page}{
  \centering
  \begin{beamercolorbox}[sep=8pt,center]{subsection title}
    \usebeamerfont{subsection title}\insertsubsection\par
  \end{beamercolorbox}
}
% Prevent slide breaks in the middle of a paragraph
\widowpenalties 1 10000
\raggedbottom
\AtBeginPart{
  \frame{\partpage}
}
\AtBeginSection{
  \ifbibliography
  \else
    \frame{\sectionpage}
  \fi
}
\AtBeginSubsection{
  \frame{\subsectionpage}
}
\usepackage{iftex}
\ifPDFTeX
  \usepackage[T1]{fontenc}
  \usepackage[utf8]{inputenc}
  \usepackage{textcomp} % provide euro and other symbols
\else % if luatex or xetex
  \usepackage{unicode-math} % this also loads fontspec
  \defaultfontfeatures{Scale=MatchLowercase}
  \defaultfontfeatures[\rmfamily]{Ligatures=TeX,Scale=1}
\fi
\usepackage{lmodern}
\ifPDFTeX\else
  % xetex/luatex font selection
\fi
% Use upquote if available, for straight quotes in verbatim environments
\IfFileExists{upquote.sty}{\usepackage{upquote}}{}
\IfFileExists{microtype.sty}{% use microtype if available
  \usepackage[]{microtype}
  \UseMicrotypeSet[protrusion]{basicmath} % disable protrusion for tt fonts
}{}
\makeatletter
\@ifundefined{KOMAClassName}{% if non-KOMA class
  \IfFileExists{parskip.sty}{%
    \usepackage{parskip}
  }{% else
    \setlength{\parindent}{0pt}
    \setlength{\parskip}{6pt plus 2pt minus 1pt}}
}{% if KOMA class
  \KOMAoptions{parskip=half}}
\makeatother
\setlength{\emergencystretch}{3em} % prevent overfull lines
\providecommand{\tightlist}{%
  \setlength{\itemsep}{0pt}\setlength{\parskip}{0pt}}
% Auto-generated by infrastructure/rendering/_pdf_latex_helpers.py::extract_math_font_preamble.
% Minimal header passed to Pandoc via -H so Beamer slides pick up the same math font
% as the combined-PDF path. Adjust the manuscript's preamble.md to change the font globally.
\usepackage{unicode-math}
\setmathfont{latinmodern-math.otf}

% Slide layout defaults for warning-clean scientific decks.
\usepackage{etoolbox}
\IfFileExists{xurl.sty}{\usepackage{xurl}}{}
\IfFileExists{seqsplit.sty}{\usepackage{seqsplit}}{\newcommand{\seqsplit}[1]{#1}}
\protected\def\breakseq#1{\seqsplit{#1}}
\protected\def\breaktt#1{\begingroup\ttfamily\seqsplit{#1}\endgroup}
\setlength{\emergencystretch}{6em}
\tolerance=5000
\hbadness=10000
\hfuzz=1pt
\setlength{\tabcolsep}{2pt}
\AtBeginEnvironment{longtable}{\tiny\renewcommand{\arraystretch}{0.86}\setlength{\tabcolsep}{1pt}}
\AtBeginEnvironment{tabular}{\tiny\renewcommand{\arraystretch}{0.86}\setlength{\tabcolsep}{1pt}}

% Natbib and cross-reference fallbacks — slides don't load natbib
% or cleveref, but combined-PDF manuscript prose may emit these
% commands. The fallback renders citations as a bracketed key list
% and cross-references as detokenized labels so slides don't fail on
% undefined control sequences or raw underscores. \providecommand is
% a no-op if packages load later.
\providecommand{\citep}[1]{[#1]}
\providecommand{\citet}[1]{#1}
\providecommand{\citealp}[1]{#1}
\providecommand{\citeauthor}[1]{#1}
\providecommand{\citeyear}[1]{#1}
\providecommand{\cref}[1]{\texttt{\detokenize{#1}}}
\providecommand{\Cref}[1]{\texttt{\detokenize{#1}}}

\newtheorem{proposition}{Proposition}
\newtheorem{hypothesis}{Hypothesis}
\usepackage{bookmark}
\IfFileExists{xurl.sty}{\usepackage{xurl}}{} % add URL line breaks if available
\urlstyle{same}
% fallback for those not using the hyperref driver hyperxmp:
\makeatletter
\@ifundefined{xmpquote}{\newcommand{\xmpquote}[1]{#1}}{}
\makeatother
\hypersetup{
  hidelinks,
  pdfcreator={LaTeX via pandoc}}

\author{\texorpdfstring{}{}}
\date{}

\begin{document}

\section{Introduction}\label{sec:introduction}

\begin{frame}[allowframebreaks]{Introduction}
An illusion stimulus is not the same thing as an illusion report. A
stimulus can be physically specified, generated, encoded, and inspected
without establishing what every observer will see or hear. This
distinction is central to reproducible cognitive science: the physical
manipulation must be auditable, while observer-level effects require a
task, calibrated presentation conditions, randomization, and data.

Visual illusion classifications are useful maps rather than exhaustive
or universally agreed ontologies \breakseq{[@gregory1997visual]}. Auditory and
audiovisual families add spectral structure, channel separation,
temporal coincidence, and spatial alignment. The Sound-Induced Flash
literature illustrates the boundary particularly clearly: the physical
contrast is a flash paired with different beep counts, while the
perceptual claim requires a listener, timing, and a response task
\breakseq{[@hirst2020sound; @shams2000sifi]}.

DuckRabbit therefore treats an illusion as a typed stimulus-construction
problem. Each catalog entry declares modality, mechanism, perceptual
signature, cognitive process, requirements, evidence status, output
kind, and implementation status. The taxonomy is orthogonal and
explicitly provisional; it is an engineering index linked to sources,
not a replacement for domain-specific theory.

The release also treats the package as research software rather than as
an unannotated collection of media files. Reproducible computational
research depends on preserving the code, inputs, environment
assumptions, and execution path needed to recreate a result
\breakseq{[@sandve2013simple; @wilson2014bestpractices]}. The FAIR literature
extends that responsibility to algorithms, tools, and workflows, while
software-citation guidance emphasizes identifying the exact software
object and version used \breakseq{[@wilkinson2016fair;
@lamprecht2020fairsoftware; @smith2016software]}. DuckRabbit applies
these principles narrowly: it makes the construction pipeline and its
generated artifacts citable and inspectable, but it does not treat
metadata quality as evidence of a perceptual effect.

The package is organized around three falsifiable engineering
hypotheses:

\begin{enumerate}
\tightlist
\item
  Identical typed requests produce identical canonical arrays and
  canonical digests.
\item
  Encoders and decoders preserve declared media facts within an explicit
  format-specific tolerance, or verification fails.
\item
  Parameter manipulations change measurable physical stimulus
  properties; any observer-level interpretation remains a preregistered
  hypothesis until data exist.
\end{enumerate}

This boundary has a nineteenth-century scientific lineage. Fechner's
\emph{Elemente der Psychophysik} formalized the problem of relating
controlled stimulus differences to measured sensation, while
Wheatstone's binocular-vision experiments made the distinction between
the physical images delivered to two eyes and the resulting depth
interpretation experimentally explicit \breakseq{[@fechner1860psychophysics;
@wheatstone1838binocular]}. Helmholtz's physiological optics then
integrated measurement of the eye, visual geometry, and theories of
perceptual inference \breakseq{[@helmholtz1867optics]}. DuckRabbit does not
reproduce those historical experiments; it inherits their methodological
lesson that stimulus conditions and observer inferences must be
represented as distinct objects.

That lineage has deeper and less geographically narrow roots. Ptolemy's
ancient \emph{Optics}, Ibn al-Haytham's medieval Arabic optics, and the
Mohist Canon's early Chinese discussions of optics and mechanics show
that questions about image formation, geometry, knowledge, and evidence
were developed across multiple intellectual traditions
\breakseq{[@ptolemy1996optics; @alhazen1989optics; @mozi2023canons;
@dai2015chinesescience]}. Early-modern discussions sharpened the
distinction between a physical presentation and an inferred percept:
Kircher documented optical display, Berkeley analyzed learned spatial
inference, and the Molyneux--Cheselden record made cross-sensory
transfer and observer testimony explicit problems \breakseq{[@kircher1646light;
@berkeley1709vision; @molyneux2020problem; @cheselden1728sight]}.
DuckRabbit cites these sources as historical and methodological
precedents, not as evidence that its generated files reproduce ancient,
early-modern, or clinical observations.

The present catalog contains 18 entries, of which 17 are implemented, 0
are planned, and 1 require external fixtures. This status vocabulary
makes breadth visible without claiming that a finite package covers all
known phenomena.
\end{frame}

\end{document}
